% 357_SM_Teilchen_GF27_Zusatz_De_ch.tex
% Johann Pascher, 7. September 2026
% Standardmodell-Teilchen aus GF(27)* — Belegstatus der Quantenzahlen (Zusatztabelle)
% Dok. 346 / 347 / 348 / 349 / 353
\providecommand{\BM}{\textbf{[B]}}
\providecommand{\KM}{\textbf{[K]}}
\providecommand{\SM}{\textbf{[S]}}
\providecommand{\XM}{\textbf{[X]}}

\phantomsection
\addcontentsline{toc}{chapter}{%
  Standardmodell-Quantenzahlen aus GF$(27)^*$ --- Belegstatus}

\section*{Statuszeichen}
\begin{center}
\begin{tabular}{@{}lp{10.5cm}@{}}
\textbf{[B]} & \emph{Belegt} --- algebraisch bewiesen. \\[3pt]
\textbf{[K]} & \emph{Konfirmiert} --- numerisch bestätigt. \\[3pt]
\textbf{[S]} & \emph{Spekulativ} --- offen, nicht abgeleitet. \\[3pt]
\textbf{[X]} & \emph{Experimentell} --- gesichert; FFGFT-Herleitung offen. \\
\end{tabular}
\end{center}

\section*{Beschreibung}

Diese Tabelle (Dok.~357) ergänzt das Teilchendiagramm (Dok.~356) und dokumentiert den Belegstatus aller Quantenzahlen des
Standardmodells, die aus der multiplikativen Gruppe GF$(27)^* =
\mathrm{GF}(3^3)^*$ algebraisch hergeleitet wurden.
Alle abgeschlossenen Herleitungen sind ohne freie Parameter.

\subsection*{Aufbau des Teilchendiagramms (Dok.~356)}

Jedes Fermion ist als konzentrisches Kreisdiagramm dargestellt.
Die Ringe von innen nach außen kodieren physikalische Eigenschaften
in der Reihenfolge ihrer algebraischen Einbettungstiefe:

\begin{center}
\begin{tabular}{@{}lp{10cm}@{}}
\textbf{Weißer Innenkreis} & Teilchensymbol und Ladung $Q$; links daneben $Y_L$ (linkshändig) und $Y_R$ (rechtshändig) \\[3pt]
\textbf{Farbring} (nur Quarks) & $r/g/b$-Sektoren = Farbladung \BM\ (Dok.~346);
  direkt nach innen, weil Gluonen die Farbe einschließen (Confinement) \\[3pt]
\textbf{Mittelring} & Links violett $= S_d =$ linkshändig \BM;
  rechts orange $= S_u =$ rechtshändig \BM\ (Dok.~349);
  grau $= \nu_R$ im SM absent \BM/\SM \\[3pt]
\textbf{Außenring} & Galois-Orbit-Farbe: Marineblau $=$ O1/O3 Gen.\,1/3 Leptonen;
  Himmelblau $=$ O3 Gen.\,2 Leptonen; Petrol $=$ O2 Gen.\,1/3 Quarks;
  Beere $=$ O4 Gen.\,2 Quarks \BM\ (Dok.~348) \\[3pt]
\textbf{Rot gestrichelt} & SU$(2)_L$-Dublett-Paar (je 2 Teilchen) \BM\ (Dok.~349) \\[3pt]
\textbf{Schwarz gestrichelt} & Generation~3 \\
\end{tabular}
\end{center}

\section{Belegstatus der Quantenzahlen}

\subsection*{Vollständig belegte Eigenschaften \BM}

\begin{description}

\item[Elektrische Ladung $Q$] (Dok.~346, 347 \BM)\\
$Q \in \{0, -1, +2/3, -1/3\}$ aus QR-Bit und $N \bmod 3$.
Alle vier Werte folgen ohne freie Parameter aus der Galois-Klassifikation
von GF$(27)^*$.

\item[Hyperladung $Y$] (Dok.~347 Satz~A \BM)\\
$Y = -1 + (4/3)\,N_\mathrm{QR}$ mit $N_\mathrm{QR}$ = Legendre-Symbol
in GF$(27)^*$. Der NQR-Bit ist exakt aus der Galois-Struktur bestimmt.

\item[Schwacher Isospin $I_3$] (Dok.~347 Satz~B \BM)\\
$I_3 = \pm 1/2$ aus dem $jE_7$-Projektor (Su/Sd-Trennung).
Die Gradparität in Cl$(6)$ erzwingt die Zuordnung.

\item[Gell-Mann-Nishijima $Q = I_3 + Y/2$] (Dok.~347 Satz~C \BM)\\
Vollständig algebraisch aus der Galois-Struktur ohne freien Parameter.

\item[Farbladung $r/g/b$] (Dok.~346 Satz~C \BM)\\
$Z_3$-Ladungsoperator $N \bmod 3$ in GF$(3)$. Drei Farben = drei
Gruppenelemente; die Dreizahl ist strukturell erzwungen.

\item[Chiralität $S_u / S_d$] (Dok.~349 \BM)\\
$S_u$ = gerade Cl$(6)$-Grade = rechtshändig;
$S_d$ = ungerade = linkshändig. Bestimmung über $\gamma_5$-Eigenwert.

\item[$SU(2)_L$-Dublett-Struktur] (Dok.~349 Satz~D \BM)\\
Das Dublett entspricht dem $S_d$-Sektor (ungerade Grade).
Das rechtshändige Neutrino $\nu_R$ ist strukturell absent.

\item[3 Generationen] (Dok.~348 Satz~A \BM)\\
Drei Frobenius-Zyklen der Nullstellen-Orbits in GF$(27)^*$.
Die Ordnung 3 ist durch die Galois-Erweiterungsstruktur erzwungen.

\item[Permutationsmatrix Quarks $\times$ Leptonen] (Dok.~348 Satz~B \BM)\\
Spur-Bilinearform GF$(27)$/GF$(3)$: $f_3 \times f_4$-Struktur.
Die Null-Kopplungen sind algebraisch bestimmt.

\item[Kopplungsgewichte $1/\sqrt{2}$] (Dok.~348 Satz~C und Dok.~353 \BM)\\
$f_1 \times f_3$-Bilinearform; zwei von drei Kopplungen aktiv.
Impliziert Koide-Amplitude $d/c = \sqrt{2}$ \BM\ (Dok.~353).

\item[Koide-Amplitude $d/c = \sqrt{2}$] (Dok.~353 \BM)\\
$Z_3$-Gleichverteilung in GF$(27)^*$ erzwingt $d/c = \sqrt{2}$,
äquivalent zur Koide-Bedingung $Q = 2/3$.
Abgeschlossen 7.~September 2026 (vormals \SM).

\item[8 Gluonen] (Dok.~339 \BM)\\
$4 \times Z_{13}$-Orbits $\times$ $2$ $Z_2$-Paritäten.
Die Gluonzahl 8 ist strukturell erzwungen.

\end{description}

\subsection*{Numerisch bestätigt \KM}

\begin{description}

\item[Elektroschwache Gruppe $Z_2$] (Dok.~336 \KM)\\
$\mathrm{Gal}(\mathrm{GF}(9)/\mathrm{GF}(3)) = Z_2$.
Numerisch bestätigt.

\end{description}

\subsection*{Offen / noch nicht geschlossen \SM}

\begin{description}

\item[Generations-Zuordnung (Reihenfolge)] (Dok.~346 \SM)\\
Welches Orbit-Element entspricht Generation 1 / 2 / 3?
Die drei Orbits sind algebraisch gleichwertig; eine Zuordnung
ist nicht aus dem Formalismus abgeleitet.

\item[CKM/PMNS-Winkel vollständig] (Dok.~348 Satz~D \SM)\\
Die Struktur der Mischungsmatrizen ist belegt (Sätze B, C \BM);
die vollständigen Winkelwerte erfordern weitere algebraische
Arbeit über die Galois-Bilinearform.

\end{description}

\subsection*{Experimentell gesichert, FFGFT-Herleitung offen \XM}

\begin{description}

\item[Higgs-Boson $H^0$] (Dok.~019 \XM)\\
$\tilde{T} \cdot m = 1$ ersetzt den Higgs-Mechanismus in FFGFT
(Dok.~019). Das Boson ist experimentell nachgewiesen;
die formale FFGFT-Herleitung des Mechanismus bleibt offen.

\end{description}
